-% Diagrama de secuencia, Análisis de informe clínico
% Uso: -% Diagrama de secuencia, Análisis de informe clínico
% Uso: -% Diagrama de secuencia, Análisis de informe clínico
% Uso: -% Diagrama de secuencia, Análisis de informe clínico
% Uso: \input{figs/seq-analisis} dentro de figure+\centering+\resizebox{\textwidth}{!}{...}
%
% Participantes (x): Usuario=0, Frontend=3.6, Backend=7.2, AI Engine=11
% Tiempo (y): crece hacia abajo (negativo en tikz)

\begin{tikzpicture}[
 every node/.style={font=\small\sffamily},
 % Cabecera de participante
 phead/.style={
 draw=aquaESI!80, fill=aquaESI!20, thick,
 rectangle, rounded corners=3pt,
 minimum width=2.4cm, minimum height=0.75cm, align=center
 },
 % Línea de vida
 life/.style={draw=gray!40, thick, dashed},
 % Mensaje síncrono (flecha llena)
 sync/.style={->, >=Stealth, thick, aquaESI!80},
 % Mensaje de retorno (flecha discontinua)
 ret/.style={->, >=Stealth, dashed, thick, gray!60},
 % Nota lateral
 note/.style={
 draw=ocre!60, fill=ocre!10,
 rectangle, rounded corners=2pt,
 font=\scriptsize\sffamily, inner sep=3pt, align=left
 },
 % Etiqueta de mensaje
 msglbl/.style={font=\scriptsize\sffamily, fill=white,
 inner sep=1pt, above, midway}
]

%% ---- COORDENADAS X de cada participante ----
\def\xu{0} % Usuario
\def\xf{3.6} % Frontend
\def\xb{7.2} % Backend
\def\xa{11.0} % AI Engine

%% ---- CABECERAS ----
\node[phead] at (\xu, 0) {\textbf{Usuario}};
\node[phead] at (\xf, 0) {\textbf{Frontend}};
\node[phead] at (\xb, 0) {\textbf{Backend}};
\node[phead] at (\xa, 0) {\textbf{AI Engine}};

%% ---- LÍNEAS DE VIDA ----
\draw[life] (\xu,-0.38) -- (\xu,-13.0);
\draw[life] (\xf,-0.38) -- (\xf,-13.0);
\draw[life] (\xb,-0.38) -- (\xb,-13.0);
\draw[life] (\xa,-0.38) -- (\xa,-13.0);

%% ---- ACTIVACIONES (barras sobre línea de vida) ----
% Backend activo durante la llamada al AI Engine
\filldraw[draw=aquaESI!70, fill=aquaESI!25]
 (\xb-0.12,-2.2) rectangle (\xb+0.12,-11.6);
% AI Engine activo durante la predicción
\filldraw[draw=turquesa!70, fill=turquesa!20]
 (\xa-0.12,-2.8) rectangle (\xa+0.12,-5.8);
% AI Engine activo durante la atribución, mucho más larga
\filldraw[draw=turquesa!70, fill=turquesa!20]
 (\xa-0.12,-8.0) rectangle (\xa+0.12,-10.8);

%% ---- MENSAJES ----

% 1. Usuario → Frontend
\draw[sync] (\xu,-1.0) -- node[msglbl]{1:\enspace envía informe clínico} (\xf,-1.0);

% 2. Frontend → Backend (WebSocket)
\draw[sync] (\xf,-1.8) -- node[msglbl]{2:\enspace [WS] analyze\_report} (\xb,-1.8);

% 3. Backend: despacha comando (auto-mensaje)
\draw[sync, rounded corners=4pt]
 (\xb,-2.4) -- ++(0.6,0) -- ++(0,-0.55) -- node[right, font=\scriptsize\sffamily]
 {3: dispatch(AnalyzeReport)} ++(-0.6,0);

% 4. Backend → AI Engine
\draw[sync] (\xb,-3.6) -- node[msglbl]{4:\enspace POST /predict + /summarize/stream (paralelo)} (\xa,-3.6);

% 5. AI Engine → Backend (respuesta)
\draw[ret] (\xa,-5.2) -- node[msglbl]{5:\enspace \{codes, confidence\_scores\}} (\xb,-5.2);

% 6. Backend: despacha evento (auto-mensaje)
\draw[sync, rounded corners=4pt]
 (\xb,-5.9) -- ++(0.6,0) -- ++(0,-0.5) -- node[right, font=\scriptsize\sffamily]
 {6: dispatch(ReceiveAIPrediction) → Event Store} ++(-0.6,0);

% 7. Backend → Frontend: tarjeta de códigos
\draw[ret] (\xb,-6.9) -- node[msglbl]{7:\enspace [WS] analysis\_card\_received: codes} (\xf,-6.9);

% 8. Backend → AI Engine: atribución en segunda llamada, no bloquea
\draw[sync] (\xb,-7.8) -- node[msglbl]{8:\enspace POST /explain (no bloquea)} (\xa,-7.8);

% 9. Backend → Frontend: resumen en streaming (token a token)
\draw[ret] (\xb,-8.7) -- node[msglbl]{9:\enspace [WS] summary\_token ($\times$N)} (\xf,-8.7);

% 10. Backend → Frontend: resumen completo
\draw[ret] (\xb,-9.5) -- node[msglbl]{10:\enspace [WS] card summary + analysis\_complete} (\xf,-9.5);

% 11. AI Engine → Backend: términos explicativos
\draw[ret] (\xa,-11.0) -- node[msglbl]{11:\enspace \{triggers, method\}} (\xb,-11.0);

% 12. Backend → Frontend: términos, cuando llegan
\draw[ret] (\xb,-11.8) -- node[msglbl]{12:\enspace [WS] triggers\_received} (\xf,-11.8);

% 13. Frontend → Usuario
\draw[ret] (\xf,-12.6) -- node[msglbl]{13:\enspace completa los términos de cada código} (\xu,-12.6);

%% ---- NOTAS LATERALES ----
\node[note, right=0.3cm] at (\xa,-4.4)
 {Inferencia BERT\\y resumen (LLM)};
\node[note, right=0.3cm] at (\xa,-9.4)
 {Atribución por\\enmascarado:\\dos órdenes de\\magnitud más lenta\\que la predicción};

\end{tikzpicture}
 dentro de figure+\centering+\resizebox{\textwidth}{!}{...}
%
% Participantes (x): Usuario=0, Frontend=3.6, Backend=7.2, AI Engine=11
% Tiempo (y): crece hacia abajo (negativo en tikz)

\begin{tikzpicture}[
 every node/.style={font=\small\sffamily},
 % Cabecera de participante
 phead/.style={
 draw=aquaESI!80, fill=aquaESI!20, thick,
 rectangle, rounded corners=3pt,
 minimum width=2.4cm, minimum height=0.75cm, align=center
 },
 % Línea de vida
 life/.style={draw=gray!40, thick, dashed},
 % Mensaje síncrono (flecha llena)
 sync/.style={->, >=Stealth, thick, aquaESI!80},
 % Mensaje de retorno (flecha discontinua)
 ret/.style={->, >=Stealth, dashed, thick, gray!60},
 % Nota lateral
 note/.style={
 draw=ocre!60, fill=ocre!10,
 rectangle, rounded corners=2pt,
 font=\scriptsize\sffamily, inner sep=3pt, align=left
 },
 % Etiqueta de mensaje
 msglbl/.style={font=\scriptsize\sffamily, fill=white,
 inner sep=1pt, above, midway}
]

%% ---- COORDENADAS X de cada participante ----
\def\xu{0} % Usuario
\def\xf{3.6} % Frontend
\def\xb{7.2} % Backend
\def\xa{11.0} % AI Engine

%% ---- CABECERAS ----
\node[phead] at (\xu, 0) {\textbf{Usuario}};
\node[phead] at (\xf, 0) {\textbf{Frontend}};
\node[phead] at (\xb, 0) {\textbf{Backend}};
\node[phead] at (\xa, 0) {\textbf{AI Engine}};

%% ---- LÍNEAS DE VIDA ----
\draw[life] (\xu,-0.38) -- (\xu,-13.0);
\draw[life] (\xf,-0.38) -- (\xf,-13.0);
\draw[life] (\xb,-0.38) -- (\xb,-13.0);
\draw[life] (\xa,-0.38) -- (\xa,-13.0);

%% ---- ACTIVACIONES (barras sobre línea de vida) ----
% Backend activo durante la llamada al AI Engine
\filldraw[draw=aquaESI!70, fill=aquaESI!25]
 (\xb-0.12,-2.2) rectangle (\xb+0.12,-11.6);
% AI Engine activo durante la predicción
\filldraw[draw=turquesa!70, fill=turquesa!20]
 (\xa-0.12,-2.8) rectangle (\xa+0.12,-5.8);
% AI Engine activo durante la atribución, mucho más larga
\filldraw[draw=turquesa!70, fill=turquesa!20]
 (\xa-0.12,-8.0) rectangle (\xa+0.12,-10.8);

%% ---- MENSAJES ----

% 1. Usuario → Frontend
\draw[sync] (\xu,-1.0) -- node[msglbl]{1:\enspace envía informe clínico} (\xf,-1.0);

% 2. Frontend → Backend (WebSocket)
\draw[sync] (\xf,-1.8) -- node[msglbl]{2:\enspace [WS] analyze\_report} (\xb,-1.8);

% 3. Backend: despacha comando (auto-mensaje)
\draw[sync, rounded corners=4pt]
 (\xb,-2.4) -- ++(0.6,0) -- ++(0,-0.55) -- node[right, font=\scriptsize\sffamily]
 {3: dispatch(AnalyzeReport)} ++(-0.6,0);

% 4. Backend → AI Engine
\draw[sync] (\xb,-3.6) -- node[msglbl]{4:\enspace POST /predict + /summarize/stream (paralelo)} (\xa,-3.6);

% 5. AI Engine → Backend (respuesta)
\draw[ret] (\xa,-5.2) -- node[msglbl]{5:\enspace \{codes, confidence\_scores\}} (\xb,-5.2);

% 6. Backend: despacha evento (auto-mensaje)
\draw[sync, rounded corners=4pt]
 (\xb,-5.9) -- ++(0.6,0) -- ++(0,-0.5) -- node[right, font=\scriptsize\sffamily]
 {6: dispatch(ReceiveAIPrediction) → Event Store} ++(-0.6,0);

% 7. Backend → Frontend: tarjeta de códigos
\draw[ret] (\xb,-6.9) -- node[msglbl]{7:\enspace [WS] analysis\_card\_received: codes} (\xf,-6.9);

% 8. Backend → AI Engine: atribución en segunda llamada, no bloquea
\draw[sync] (\xb,-7.8) -- node[msglbl]{8:\enspace POST /explain (no bloquea)} (\xa,-7.8);

% 9. Backend → Frontend: resumen en streaming (token a token)
\draw[ret] (\xb,-8.7) -- node[msglbl]{9:\enspace [WS] summary\_token ($\times$N)} (\xf,-8.7);

% 10. Backend → Frontend: resumen completo
\draw[ret] (\xb,-9.5) -- node[msglbl]{10:\enspace [WS] card summary + analysis\_complete} (\xf,-9.5);

% 11. AI Engine → Backend: términos explicativos
\draw[ret] (\xa,-11.0) -- node[msglbl]{11:\enspace \{triggers, method\}} (\xb,-11.0);

% 12. Backend → Frontend: términos, cuando llegan
\draw[ret] (\xb,-11.8) -- node[msglbl]{12:\enspace [WS] triggers\_received} (\xf,-11.8);

% 13. Frontend → Usuario
\draw[ret] (\xf,-12.6) -- node[msglbl]{13:\enspace completa los términos de cada código} (\xu,-12.6);

%% ---- NOTAS LATERALES ----
\node[note, right=0.3cm] at (\xa,-4.4)
 {Inferencia BERT\\y resumen (LLM)};
\node[note, right=0.3cm] at (\xa,-9.4)
 {Atribución por\\enmascarado:\\dos órdenes de\\magnitud más lenta\\que la predicción};

\end{tikzpicture}
 dentro de figure+\centering+\resizebox{\textwidth}{!}{...}
%
% Participantes (x): Usuario=0, Frontend=3.6, Backend=7.2, AI Engine=11
% Tiempo (y): crece hacia abajo (negativo en tikz)

\begin{tikzpicture}[
 every node/.style={font=\small\sffamily},
 % Cabecera de participante
 phead/.style={
 draw=aquaESI!80, fill=aquaESI!20, thick,
 rectangle, rounded corners=3pt,
 minimum width=2.4cm, minimum height=0.75cm, align=center
 },
 % Línea de vida
 life/.style={draw=gray!40, thick, dashed},
 % Mensaje síncrono (flecha llena)
 sync/.style={->, >=Stealth, thick, aquaESI!80},
 % Mensaje de retorno (flecha discontinua)
 ret/.style={->, >=Stealth, dashed, thick, gray!60},
 % Nota lateral
 note/.style={
 draw=ocre!60, fill=ocre!10,
 rectangle, rounded corners=2pt,
 font=\scriptsize\sffamily, inner sep=3pt, align=left
 },
 % Etiqueta de mensaje
 msglbl/.style={font=\scriptsize\sffamily, fill=white,
 inner sep=1pt, above, midway}
]

%% ---- COORDENADAS X de cada participante ----
\def\xu{0} % Usuario
\def\xf{3.6} % Frontend
\def\xb{7.2} % Backend
\def\xa{11.0} % AI Engine

%% ---- CABECERAS ----
\node[phead] at (\xu, 0) {\textbf{Usuario}};
\node[phead] at (\xf, 0) {\textbf{Frontend}};
\node[phead] at (\xb, 0) {\textbf{Backend}};
\node[phead] at (\xa, 0) {\textbf{AI Engine}};

%% ---- LÍNEAS DE VIDA ----
\draw[life] (\xu,-0.38) -- (\xu,-13.0);
\draw[life] (\xf,-0.38) -- (\xf,-13.0);
\draw[life] (\xb,-0.38) -- (\xb,-13.0);
\draw[life] (\xa,-0.38) -- (\xa,-13.0);

%% ---- ACTIVACIONES (barras sobre línea de vida) ----
% Backend activo durante la llamada al AI Engine
\filldraw[draw=aquaESI!70, fill=aquaESI!25]
 (\xb-0.12,-2.2) rectangle (\xb+0.12,-11.6);
% AI Engine activo durante la predicción
\filldraw[draw=turquesa!70, fill=turquesa!20]
 (\xa-0.12,-2.8) rectangle (\xa+0.12,-5.8);
% AI Engine activo durante la atribución, mucho más larga
\filldraw[draw=turquesa!70, fill=turquesa!20]
 (\xa-0.12,-8.0) rectangle (\xa+0.12,-10.8);

%% ---- MENSAJES ----

% 1. Usuario → Frontend
\draw[sync] (\xu,-1.0) -- node[msglbl]{1:\enspace envía informe clínico} (\xf,-1.0);

% 2. Frontend → Backend (WebSocket)
\draw[sync] (\xf,-1.8) -- node[msglbl]{2:\enspace [WS] analyze\_report} (\xb,-1.8);

% 3. Backend: despacha comando (auto-mensaje)
\draw[sync, rounded corners=4pt]
 (\xb,-2.4) -- ++(0.6,0) -- ++(0,-0.55) -- node[right, font=\scriptsize\sffamily]
 {3: dispatch(AnalyzeReport)} ++(-0.6,0);

% 4. Backend → AI Engine
\draw[sync] (\xb,-3.6) -- node[msglbl]{4:\enspace POST /predict + /summarize/stream (paralelo)} (\xa,-3.6);

% 5. AI Engine → Backend (respuesta)
\draw[ret] (\xa,-5.2) -- node[msglbl]{5:\enspace \{codes, confidence\_scores\}} (\xb,-5.2);

% 6. Backend: despacha evento (auto-mensaje)
\draw[sync, rounded corners=4pt]
 (\xb,-5.9) -- ++(0.6,0) -- ++(0,-0.5) -- node[right, font=\scriptsize\sffamily]
 {6: dispatch(ReceiveAIPrediction) → Event Store} ++(-0.6,0);

% 7. Backend → Frontend: tarjeta de códigos
\draw[ret] (\xb,-6.9) -- node[msglbl]{7:\enspace [WS] analysis\_card\_received: codes} (\xf,-6.9);

% 8. Backend → AI Engine: atribución en segunda llamada, no bloquea
\draw[sync] (\xb,-7.8) -- node[msglbl]{8:\enspace POST /explain (no bloquea)} (\xa,-7.8);

% 9. Backend → Frontend: resumen en streaming (token a token)
\draw[ret] (\xb,-8.7) -- node[msglbl]{9:\enspace [WS] summary\_token ($\times$N)} (\xf,-8.7);

% 10. Backend → Frontend: resumen completo
\draw[ret] (\xb,-9.5) -- node[msglbl]{10:\enspace [WS] card summary + analysis\_complete} (\xf,-9.5);

% 11. AI Engine → Backend: términos explicativos
\draw[ret] (\xa,-11.0) -- node[msglbl]{11:\enspace \{triggers, method\}} (\xb,-11.0);

% 12. Backend → Frontend: términos, cuando llegan
\draw[ret] (\xb,-11.8) -- node[msglbl]{12:\enspace [WS] triggers\_received} (\xf,-11.8);

% 13. Frontend → Usuario
\draw[ret] (\xf,-12.6) -- node[msglbl]{13:\enspace completa los términos de cada código} (\xu,-12.6);

%% ---- NOTAS LATERALES ----
\node[note, right=0.3cm] at (\xa,-4.4)
 {Inferencia BERT\\y resumen (LLM)};
\node[note, right=0.3cm] at (\xa,-9.4)
 {Atribución por\\enmascarado:\\dos órdenes de\\magnitud más lenta\\que la predicción};

\end{tikzpicture}
 dentro de figure+\centering+\resizebox{\textwidth}{!}{...}
%
% Participantes (x): Usuario=0, Frontend=3.6, Backend=7.2, AI Engine=11
% Tiempo (y): crece hacia abajo (negativo en tikz)

\begin{tikzpicture}[
 every node/.style={font=\small\sffamily},
 % Cabecera de participante
 phead/.style={
 draw=aquaESI!80, fill=aquaESI!20, thick,
 rectangle, rounded corners=3pt,
 minimum width=2.4cm, minimum height=0.75cm, align=center
 },
 % Línea de vida
 life/.style={draw=gray!40, thick, dashed},
 % Mensaje síncrono (flecha llena)
 sync/.style={->, >=Stealth, thick, aquaESI!80},
 % Mensaje de retorno (flecha discontinua)
 ret/.style={->, >=Stealth, dashed, thick, gray!60},
 % Nota lateral
 note/.style={
 draw=ocre!60, fill=ocre!10,
 rectangle, rounded corners=2pt,
 font=\scriptsize\sffamily, inner sep=3pt, align=left
 },
 % Etiqueta de mensaje
 msglbl/.style={font=\scriptsize\sffamily, fill=white,
 inner sep=1pt, above, midway}
]

%% ---- COORDENADAS X de cada participante ----
\def\xu{0} % Usuario
\def\xf{3.6} % Frontend
\def\xb{7.2} % Backend
\def\xa{11.0} % AI Engine

%% ---- CABECERAS ----
\node[phead] at (\xu, 0) {\textbf{Usuario}};
\node[phead] at (\xf, 0) {\textbf{Frontend}};
\node[phead] at (\xb, 0) {\textbf{Backend}};
\node[phead] at (\xa, 0) {\textbf{AI Engine}};

%% ---- LÍNEAS DE VIDA ----
\draw[life] (\xu,-0.38) -- (\xu,-13.0);
\draw[life] (\xf,-0.38) -- (\xf,-13.0);
\draw[life] (\xb,-0.38) -- (\xb,-13.0);
\draw[life] (\xa,-0.38) -- (\xa,-13.0);

%% ---- ACTIVACIONES (barras sobre línea de vida) ----
% Backend activo durante la llamada al AI Engine
\filldraw[draw=aquaESI!70, fill=aquaESI!25]
 (\xb-0.12,-2.2) rectangle (\xb+0.12,-11.6);
% AI Engine activo durante la predicción
\filldraw[draw=turquesa!70, fill=turquesa!20]
 (\xa-0.12,-2.8) rectangle (\xa+0.12,-5.8);
% AI Engine activo durante la atribución, mucho más larga
\filldraw[draw=turquesa!70, fill=turquesa!20]
 (\xa-0.12,-8.0) rectangle (\xa+0.12,-10.8);

%% ---- MENSAJES ----

% 1. Usuario → Frontend
\draw[sync] (\xu,-1.0) -- node[msglbl]{1:\enspace envía informe clínico} (\xf,-1.0);

% 2. Frontend → Backend (WebSocket)
\draw[sync] (\xf,-1.8) -- node[msglbl]{2:\enspace [WS] analyze\_report} (\xb,-1.8);

% 3. Backend: despacha comando (auto-mensaje)
\draw[sync, rounded corners=4pt]
 (\xb,-2.4) -- ++(0.6,0) -- ++(0,-0.55) -- node[right, font=\scriptsize\sffamily]
 {3: dispatch(AnalyzeReport)} ++(-0.6,0);

% 4. Backend → AI Engine
\draw[sync] (\xb,-3.6) -- node[msglbl]{4:\enspace POST /predict + /summarize/stream (paralelo)} (\xa,-3.6);

% 5. AI Engine → Backend (respuesta)
\draw[ret] (\xa,-5.2) -- node[msglbl]{5:\enspace \{codes, confidence\_scores\}} (\xb,-5.2);

% 6. Backend: despacha evento (auto-mensaje)
\draw[sync, rounded corners=4pt]
 (\xb,-5.9) -- ++(0.6,0) -- ++(0,-0.5) -- node[right, font=\scriptsize\sffamily]
 {6: dispatch(ReceiveAIPrediction) → Event Store} ++(-0.6,0);

% 7. Backend → Frontend: tarjeta de códigos
\draw[ret] (\xb,-6.9) -- node[msglbl]{7:\enspace [WS] analysis\_card\_received: codes} (\xf,-6.9);

% 8. Backend → AI Engine: atribución en segunda llamada, no bloquea
\draw[sync] (\xb,-7.8) -- node[msglbl]{8:\enspace POST /explain (no bloquea)} (\xa,-7.8);

% 9. Backend → Frontend: resumen en streaming (token a token)
\draw[ret] (\xb,-8.7) -- node[msglbl]{9:\enspace [WS] summary\_token ($\times$N)} (\xf,-8.7);

% 10. Backend → Frontend: resumen completo
\draw[ret] (\xb,-9.5) -- node[msglbl]{10:\enspace [WS] card summary + analysis\_complete} (\xf,-9.5);

% 11. AI Engine → Backend: términos explicativos
\draw[ret] (\xa,-11.0) -- node[msglbl]{11:\enspace \{triggers, method\}} (\xb,-11.0);

% 12. Backend → Frontend: términos, cuando llegan
\draw[ret] (\xb,-11.8) -- node[msglbl]{12:\enspace [WS] triggers\_received} (\xf,-11.8);

% 13. Frontend → Usuario
\draw[ret] (\xf,-12.6) -- node[msglbl]{13:\enspace completa los términos de cada código} (\xu,-12.6);

%% ---- NOTAS LATERALES ----
\node[note, right=0.3cm] at (\xa,-4.4)
 {Inferencia BERT\\y resumen (LLM)};
\node[note, right=0.3cm] at (\xa,-9.4)
 {Atribución por\\enmascarado:\\dos órdenes de\\magnitud más lenta\\que la predicción};

\end{tikzpicture}
